% DesignFix ASE 2026 ACM review-format skeleton
% Keep the existing figure files in the Overleaf folder:
%   images/DesignFix_System_Diagram_v4.png
%   images/1_anno.png
%   images/2_anno_alt.png
%   images/codefix_passrate.png
%   images/explanation_passrate.png
%   images/surrounding_code_passrate.png
%   images/Task_Completion_Chart_boxplot.png
%   images/NASA_TLX_enlarged.png

\begin{filecontents*}[overwrite]{references.bib}
@inproceedings{mehrpour2019active,
  author    = {Mehrpour, Sahar and LaToza, Thomas D. and Kindi, Rahul K.},
  title     = {Active Documentation: Helping Developers Follow Design Decisions},
  booktitle = {Proceedings of the IEEE Symposium on Visual Languages and Human-Centric Computing},
  pages     = {87--96},
  year      = {2019}
}

@inproceedings{latoza2014microtask,
  author    = {LaToza, Thomas D. and Towne, W. Ben and Adriano, Christian M. and van der Hoek, Andr{\'e}},
  title     = {Microtask Programming: Building Software with a Crowd},
  booktitle = {Proceedings of the 27th Annual ACM Symposium on User Interface Software and Technology},
  pages     = {43--54},
  year      = {2014},
  doi       = {10.1145/2642918.2647349}
}

@article{latoza2019microtaskTSE,
  author  = {LaToza, Thomas D. and Di Lecce, Arturo and Ricci, Fabio and Towne, W. Ben and van der Hoek, Andr{\'e}},
  title   = {Microtask Programming},
  journal = {IEEE Transactions on Software Engineering},
  volume  = {45},
  number  = {11},
  pages   = {1106--1124},
  year    = {2019},
  doi     = {10.1109/TSE.2018.2823327}
}

@inproceedings{mehrpour2020rulepad,
  author    = {Mehrpour, Sahar and LaToza, Thomas D. and Sarvari, Hamed},
  title     = {RulePad: Interactive Authoring of Checkable Design Rules},
  booktitle = {Proceedings of the 28th ACM Joint Meeting on European Software Engineering Conference and Symposium on the Foundations of Software Engineering},
  pages     = {386--397},
  year      = {2020},
  doi       = {10.1145/3368089.3409751}
}

@article{mcdonald2019reliability,
  author  = {McDonald, Nora and Schoenebeck, Sarita and Forte, Andrea},
  title   = {Reliability and Inter-rater Reliability in Qualitative Research: Norms and Guidelines for {CSCW} and {HCI} Practice},
  journal = {Proceedings of the ACM on Human-Computer Interaction},
  volume  = {3},
  number  = {CSCW},
  pages   = {72:1--72:23},
  year    = {2019},
  doi     = {10.1145/3359174}
}


@article{mehrpour2023canstatic,
  author  = {Mehrpour, Sahar and LaToza, Thomas D.},
  title   = {Can Static Analysis Tools Find More Defects? A Qualitative Study of Design Rule Violations Found by Code Review},
  journal = {Empirical Software Engineering},
  volume  = {28},
  number  = {1},
  pages   = {5},
  year    = {2023}
}


@misc{peng2023impact,
  author       = {Peng, Sida and Kalliamvakou, Eirini and Cihon, Peter and Demirer, Mert},
  title        = {The Impact of AI on Developer Productivity: Evidence from GitHub Copilot},
  howpublished = {arXiv preprint arXiv:2302.06590},
  year         = {2023}
}

@inproceedings{liang2024survey,
  author    = {Liang, Jenny T. and Yang, Chenyang and Myers, Brad A.},
  title     = {A Large-Scale Survey on the Usability of AI Programming Assistants: Successes and Challenges},
  booktitle = {Proceedings of the 46th IEEE/ACM International Conference on Software Engineering},
  year      = {2024}
}

@inproceedings{johnson2023trust,
  author    = {Johnson, Brittany and Bird, Christian and Ford, Denae and Forsgren, Nicole and Zimmermann, Thomas},
  title     = {Make Your Tools Sparkle with Trust: The PICSE Framework for Trust in Software Tools},
  booktitle = {Proceedings of the IEEE/ACM 45th International Conference on Software Engineering: Software Engineering in Practice},
  pages     = {409--419},
  year      = {2023}
}

@article{russo2024navigating,
  author  = {Russo, Daniel},
  title   = {Navigating the Complexity of Generative AI Adoption in Software Engineering},
  journal = {ACM Transactions on Software Engineering and Methodology},
  volume  = {33},
  number  = {5},
  pages   = {1--50},
  year    = {2024}
}

@inproceedings{alhaque2025evolution,
  author    = {Al Haque, Ehsan and Brown, Chris and LaToza, Thomas D. and Johnson, Brittany},
  title     = {The Evolution of Information Seeking in Software Development: Understanding the Role and Impact of AI Assistants},
  booktitle = {Proceedings of the ACM International Conference on the Foundations of Software Engineering Companion},
  year      = {2025}
}

@inproceedings{nachtigall2022usability,
  author    = {Nachtigall, Marcus and Schlichtig, Michael and Bodden, Eric},
  title     = {A Large-Scale Study of Usability Criteria Addressed by Static Analysis Tools},
  booktitle = {Proceedings of the ACM SIGSOFT International Symposium on Software Testing and Analysis},
  pages     = {532--543},
  year      = {2022}
}

@article{guo2023falsepositive,
  author  = {Guo, Zhaoqiang and Tan, Tingting and Liu, Shiran and others},
  title   = {Mitigating False Positive Static Analysis Warnings: Progress, Challenges, and Opportunities},
  journal = {IEEE Transactions on Software Engineering},
  volume  = {49},
  number  = {12},
  pages   = {5154--5188},
  year    = {2023}
}

@article{li2022erosion,
  author  = {Li, Ruiyin and Liang, Peng and Soliman, Mohamed and Avgeriou, Paris},
  title   = {Understanding Software Architecture Erosion: A Systematic Mapping Study},
  journal = {Journal of Software: Evolution and Process},
  volume  = {34},
  number  = {3},
  pages   = {e2423},
  year    = {2022}
}

@article{liu2023deepsmell,
  author  = {Liu, Hui and Jin, Jiahao and Xu, Zhifeng and Zou, Yanzhen and Bu, Yifan and Zhang, Lu},
  title   = {Deep Learning Based Code Smell Detection},
  journal = {IEEE Transactions on Software Engineering},
  volume  = {47},
  number  = {9},
  pages   = {1811--1837},
  year    = {2021}
}

@article{etemadi2022sorald,
  author  = {Etemadi, Khashayar and Harrand, Nicolas and Larsen, Simon and others},
  title   = {Sorald: Automatic Patch Suggestions for SonarQube Static Analysis Violations},
  journal = {IEEE Transactions on Dependable and Secure Computing},
  volume  = {20},
  number  = {4},
  pages   = {2794--2810},
  year    = {2023}
}

@article{wadhwa2024core,
  author  = {Wadhwa, Nalin and Pradhan, Jui and Sonwane, Atharv and others},
  title   = {CORE: Resolving Code Quality Issues Using LLMs},
  journal = {Proceedings of the ACM on Software Engineering},
  volume  = {1},
  number  = {FSE},
  year    = {2024}
}

@inproceedings{xia2022alpharepair,
  author    = {Xia, Chunqiu Steven and Zhang, Lingming},
  title     = {Less Training, More Repairing Please: Revisiting Automated Program Repair via Zero-Shot Learning},
  booktitle = {Proceedings of the ACM Joint European Software Engineering Conference and Symposium on the Foundations of Software Engineering},
  pages     = {959--971},
  year      = {2022}
}

@inproceedings{xia2023era,
  author    = {Xia, Chunqiu Steven and Wei, Yuxiang and Zhang, Lingming},
  title     = {Automated Program Repair in the Era of Large Pre-Trained Language Models},
  booktitle = {Proceedings of the 45th IEEE/ACM International Conference on Software Engineering},
  pages     = {1482--1494},
  year      = {2023}
}

@inproceedings{xia2024chatrepair,
  author    = {Xia, Chunqiu Steven and Zhang, Lingming},
  title     = {Automated Program Repair via Conversation: Fixing 162 out of 337 Bugs for \$0.42 Each Using ChatGPT},
  booktitle = {Proceedings of the ACM SIGSOFT International Symposium on Software Testing and Analysis},
  pages     = {819--831},
  year      = {2024}
}

@misc{xiang2024srepair,
  author       = {Xiang, Jiahong and Xu, Xiaoyang and Kong, Fanchu and others},
  title        = {How Far Can We Go with Practical Function-Level Program Repair?},
  howpublished = {arXiv preprint arXiv:2404.12833},
  year         = {2024}
}

@inproceedings{bouzenia2025repairagent,
  author    = {Bouzenia, Islem and Devanbu, Premkumar and Pradel, Michael},
  title     = {RepairAgent: An Autonomous, LLM-Based Agent for Program Repair},
  booktitle = {Proceedings of the 47th IEEE/ACM International Conference on Software Engineering},
  year      = {2025}
}

@inproceedings{joshi2023ring,
  author    = {Joshi, Harshit and Cambronero, Jos{\'e} and Gulwani, Sumit and Le, Vu and Radi{\v{c}}ek, Ivan and Verbruggen, Gust},
  title     = {Repair Is Nearly Generation: Multilingual Program Repair with LLMs},
  booktitle = {Proceedings of the AAAI Conference on Artificial Intelligence},
  volume    = {37},
  pages     = {5131--5140},
  year      = {2023}
}

@inproceedings{wang2023rapgen,
  author    = {Wang, Weishi and Wang, Yue and Joty, Shafiq and Hoi, Steven C. H.},
  title     = {RAP-Gen: Retrieval-Augmented Patch Generation with CodeT5 for Automatic Program Repair},
  booktitle = {Proceedings of the ACM Joint European Software Engineering Conference and Symposium on the Foundations of Software Engineering},
  pages     = {146--158},
  year      = {2023}
}

@inproceedings{jin2023inferfix,
  author    = {Jin, Matthew and Shahriar, Syed and Tufano, Michele and others},
  title     = {InferFix: End-to-End Program Repair with LLMs over Retrieval-Augmented Prompts},
  booktitle = {Proceedings of the ESEC/FSE Industry Track},
  pages     = {1646--1656},
  year      = {2023}
}

@article{barke2023grounded,
  author  = {Barke, Shraddha and James, Michael B. and Polikarpova, Nadia},
  title   = {Grounded Copilot: How Programmers Interact with Code-Generating Models},
  journal = {Proceedings of the ACM on Programming Languages},
  volume  = {7},
  number  = {OOPSLA1},
  pages   = {85--111},
  year    = {2023}
}

@inproceedings{vaithilingam2023intellicode,
  author    = {Vaithilingam, Priyan and Glassman, Elena L. and Groenwegen, Peter and others},
  title     = {Towards More Effective AI-Assisted Programming: A Systematic Design Exploration to Improve Visual Studio IntelliCode's User Experience},
  booktitle = {Proceedings of the IEEE/ACM 45th International Conference on Software Engineering: Software Engineering in Practice},
  pages     = {185--195},
  year      = {2023}
}

@inproceedings{perry2023insecure,
  author    = {Perry, Neil and Srivastava, Megha and Kumar, Deepak and Boneh, Dan},
  title     = {Do Users Write More Insecure Code with AI Assistants?},
  booktitle = {Proceedings of the ACM SIGSAC Conference on Computer and Communications Security},
  pages     = {2785--2799},
  year      = {2023}
}

@article{li2024scot,
  author  = {Li, Jia and Li, Ge and Li, Yongmin and Jin, Zhi},
  title   = {Structured Chain-of-Thought Prompting for Code Generation},
  journal = {ACM Transactions on Software Engineering and Methodology},
  volume  = {34},
  number  = {2},
  pages   = {1--23},
  year    = {2025}
}

@article{jiang2024selfplanning,
  author  = {Jiang, Xue and Dong, Yihong and Wang, Lecheng and others},
  title   = {Self-Planning Code Generation with Large Language Models},
  journal = {ACM Transactions on Software Engineering and Methodology},
  volume  = {33},
  number  = {7},
  pages   = {1--30},
  year    = {2024}
}

@inproceedings{chen2024selfdebug,
  author    = {Chen, Xinyun and Lin, Maxwell and Sch{\"a}rli, Nathanael and Zhou, Denny},
  title     = {Teaching Large Language Models to Self-Debug},
  booktitle = {Proceedings of the International Conference on Learning Representations},
  year      = {2024}
}

@inproceedings{madaan2023selfrefine,
  author    = {Madaan, Aman and Tandon, Niket and Gupta, Prakhar and others},
  title     = {Self-Refine: Iterative Refinement with Self-Feedback},
  booktitle = {Advances in Neural Information Processing Systems},
  volume    = {36},
  year      = {2023}
}

@inproceedings{zhou2023docprompting,
  author    = {Zhou, Shuyan and Alon, Uri and Xu, Frank F. and others},
  title     = {DocPrompting: Generating Code by Retrieving the Docs},
  booktitle = {Proceedings of the International Conference on Learning Representations},
  year      = {2023}
}

@inproceedings{nashid2023cedar,
  author    = {Nashid, Noor and Sintaha, Mifta and Mesbah, Ali},
  title     = {Retrieval-Based Prompt Selection for Code-Related Few-Shot Learning},
  booktitle = {Proceedings of the 45th IEEE/ACM International Conference on Software Engineering},
  pages     = {2456--2468},
  year      = {2023}
}

@incollection{hart1988development,
  author    = {Hart, Sandra G. and Staveland, Lowell E.},
  title     = {Development of NASA-TLX (Task Load Index): Results of Empirical and Theoretical Research},
  booktitle = {Advances in Psychology},
  volume    = {52},
  pages     = {139--183},
  year      = {1988}
}

@inproceedings{collard2013srcml,
  author    = {Collard, Michael L. and Decker, Michael J. and Maletic, Jonathan I.},
  title     = {srcML: An Infrastructure for the Exploration, Analysis, and Manipulation of Source Code: A Tool Demonstration},
  booktitle = {Proceedings of the 29th IEEE International Conference on Software Maintenance},
  pages     = {516--519},
  year      = {2013}
}
\end{filecontents*}

\documentclass[sigconf,review,anonymous]{acmart}

\usepackage{graphicx}
\usepackage{xcolor}
\usepackage{booktabs}
\usepackage{listings}
\usepackage{pifont}
\usepackage{tikz}

\graphicspath{{images/}}

\newcommand{\circmark}[1]{%
  \tikz[baseline=(char.base)]{
    \node[
      shape=circle,
      draw=black!40,
      fill=black!10,
      inner sep=1.2pt,
      font=\small\bfseries
    ] (char) {#1};
  }%
}


\newcommand{\pquote}[3]{%
\vspace{2pt}
\begin{quote}\small
\textbf{#1} (\emph{#2}): \emph{``#3''}
\end{quote}
\vspace{2pt}
}

\lstset{
  basicstyle=\ttfamily\small,
  breaklines=true,
  frame=single,
  numbers=none,
  tabsize=2,
  captionpos=b,
  abovecaptionskip=5pt,
}

\begin{document}

\title{DesignFix: Context-Augmented Repair of Design Rule Violations}

%% Real author info stays in the source; the [anonymous] class option
%% prints "Anonymous Author(s)" in the review PDF. Remove the review and
%% anonymous options for the camera-ready and these will print.
\author{Mainul Hossain}
\email{mhossa6@gmu.edu}
\affiliation{%
  \institution{George Mason University}
  \city{Fairfax}
  \state{Virginia}
  \country{USA}
}

\author{Thomas D. LaToza}
\email{tlatoza@gmu.edu}
\affiliation{%
  \institution{George Mason University}
  \city{Fairfax}
  \state{Virginia}
  \country{USA}
}


\begin{comment}
\begin{abstract}
% Add the abstract here.
Design rules capture constraints that developers must follow to ensure code remains consistent with design decisions. While tools exist to detect design rule violations, understanding why code violates a rule and determining how to fix it remains challenging, particularly in unfamiliar codebases. We propose DesignFix, an approach that applies large language models to assisting developers in interpreting and repairing design rule violations. DesignFix builds on automated design rule checking to extract violations along with contextual information from the codebase including rule descriptions, rule-satisfying examples from the same project, and surrounding file context. This context augmentation grounds the LLM's reasoning in project-specific patterns rather than relying on generic code knowledge. A two-step co-generative prompting strategy separates free-form reasoning from structured output generation, improving both fix quality and explanation accuracy. Fixes are presented through an interactive diff-based review workflow enabling developers to verify and accept or reject suggestions. Through a systematic evaluation of seven prompt variants across 9,450 model runs, we found that rule-following examples and a two step approach achieved the highest performance, with 90\% code fix success, 89\% explanation correctness, and 90\% surrounding code preservation. A within-subjects user study with 14 participants demonstrated that DesignFix significantly reduced task completion time and improved correctness, while significantly lowering cognitive workload across 5 NASA-TLX dimensions. Qualitative analysis revealed that DesignFix's primary value derives from precise codebase navigation rather than code generation, with participants spending 75\% less time locating relevant code, and that its verification affordances enabled developers to detect and recover from incorrect suggestions.
\end{abstract}
\end{comment}
\begin{abstract}
Design rules encode project-specific constraints that preserve
architectural intent. Although existing checkers can detect violations,
developers must still interpret the rule, locate the repair site, and
implement a project-consistent fix. We present DesignFix, an
IDE-integrated repair workflow that augments checker-reported
violations with rule descriptions, same-project satisfying examples,
focused structural context, and source-file context. A two-step LLM
pipeline separates repair reasoning from structured extraction, then
presents the explanation and candidate patch for developer review. We
evaluated seven cumulative prompt configurations in 9,450 repair
trials. The strongest tested workflow achieved up to 90\% repair
success, 89\% explanation correctness, and 90\% surrounding-code
preservation. In a within-subjects study with 14 participants,
DesignFix reduced mean task time by 36\%, increased correct repairs
from 43\% to 93\%, and significantly reduced five of six NASA-TLX
dimensions. Participants particularly valued repair-site navigation,
with mean time to first edit decreasing from 4.9 to 2.1 minutes, while
explanation and review support helped them identify and recover from
imperfect suggestions.
\end{abstract}

%% CCS concepts are mandatory for ACM papers over two pages; concept IDs
%% verified at https://dl.acm.org/ccs (regenerate there to change them).
\begin{CCSXML}
<ccs2012>
   <concept>
       <concept_id>10011007.10011006.10011073</concept_id>
       <concept_desc>Software and its engineering~Software maintenance tools</concept_desc>
       <concept_significance>500</concept_significance>
   </concept>
   <concept>
       <concept_id>10003120.10003121.10011748</concept_id>
       <concept_desc>Human-centered computing~Empirical studies in HCI</concept_desc>
       <concept_significance>300</concept_significance>
   </concept>
</ccs2012>
\end{CCSXML}

\ccsdesc[500]{Software and its engineering~Software maintenance tools}
\ccsdesc[300]{Human-centered computing~Empirical studies in HCI}

\keywords{design rules, large language models, software maintenance, code repair, developer tools}

\maketitle

\section{Introduction}
\label{sec:introduction}
% Add section text here.
Project-specific design decisions constrain how software should be
implemented. Design rules make these constraints checkable, helping
teams preserve architectural intent and prevent defects as systems
evolve~\cite{mehrpour2019active}.

Design rule violations are both common and consequential. In an
analysis of 1,323 code-review defects from 493 pull requests, Mehrpour
and LaToza found that up to 76\% could in principle be detected
statically~\cite{mehrpour2023canstatic}. AST patterns, the rule form
studied here, covered 25\% of all defects, more than any other analysis
type, and most resulting rules were project-specific. Yet detection is
only the first step. Developers must interpret the rule, identify the
intended project-local pattern, and locate the program element that
must change. This is especially costly in unfamiliar codebases, where
design rationale and valid examples may be dispersed across
documentation and source files.
\begin{comment}
The problem is compounded by the nature of modern software development. Developers frequently join new teams, inherit legacy codebases, or contribute to open-source projects where they lack deep familiarity with the design decisions that shaped the code. In these situations, fixing a design rule violation requires not just identifying what is wrong, but understanding the underlying design intent and how to align the fix with existing patterns in the codebase. This process can involve examining multiple examples of code that satisfies the rule, tracing through class hierarchies, or understanding framework-specific conventions that may not be immediately obvious.
\end{comment}

LLMs create an opportunity to transform checker output into actionable,
project-grounded repair guidance. The central challenge is not simply
generating code, but selecting context that exposes local design intent
and presenting the result in a form developers can verify.

We present DesignFix, an IDE-integrated repair workflow that transforms
a detected design rule violation into a project-grounded explanation
and reviewable patch. DesignFix combines checker-derived rule
information, same-project examples, focused code context, and a
two-step prompting pipeline to help developers identify the repair
site, understand the intended design, and implement a codebase-consistent
fix.

We evaluate DesignFix through a systematic prompt study and a
within-subjects developer study addressing three questions:

\begin{itemize}
\item \textbf{RQ1:} How proficient are large language models in interpreting and responding to different types of coding prompts, and what limitations do they exhibit?
\item \textbf{RQ2:} How effective are LLMs in suggesting accurate code modifications and providing comprehensible explanations across diverse codebases, and what factors influence their performance?
\item \textbf{RQ3:} How does DesignFix affect developer  performance fixing design rule violations in an unfamiliar codebase?
\end{itemize}


\begin{comment}
To answer these questions, we conducted two complementary evaluations. First, we systematically explored different prompting strategies with two versions of gpt-4o, iteratively refining prompts to understand what context and formatting produces the most accurate and useful responses. We evaluated seven prompt variants across 45 design rule violations from two distinct codebases, measuring code fix success rate, Explanation Pass Rate, and Surrounding Code Pass Rate. In total, this evaluation involved 9,450 individual model runs across different prompt types, model versions, and temperature settings.

Second, we conducted a within-subjects user study with 14 participants who had varying levels of software engineering experience, ranging from 0 to 12 years. Each participant fixed one design rule violation without LLM assistance and another with DesignFix, allowing us to directly compare task performance, cognitive workload, and subjective experience.
\end{comment}

\begin{comment}
Our results show that: (1) a two-step co-generative prompting approach significantly outperforms single-shot prompts, achieving 90\% code fix success and 89\% explanation correctness; (2) participants using DesignFix completed tasks significantly faster (10.0 vs. 15.71 minutes, representing a 36\% time reduction) with substantially higher correctness (93\% vs. 43\%); and (3) DesignFix reduced cognitive workload across all NASA-TLX dimensions, with particularly large reductions in mental demand, effort, and frustration.
\end{comment}

The contributions of this paper are:

\begin{itemize}
\item \textbf{DesignFix}, a project-grounded repair workflow that
connects checker-reported violations to same-project examples, focused
context, explanations, and developer-reviewable patches.

\item \textbf{A systematic prompt evaluation} revealing how
project-specific examples, focused context, and two-step generation
shape repair and explanation quality.

\item \textbf{A controlled developer study} demonstrating improvements
in task efficiency, repair correctness, and workload while revealing
how navigation and verification support developer decision making.
\end{itemize}

\section{Background: Design Rules}
\label{sec:background}
% Add section text here.
A design decision captures an implementation choice and its rationale.
A design rule makes the resulting project-specific constraint
executable against source code, enabling tools to check whether an
implementation remains consistent with its intended design
~\cite{mehrpour2019active}.

\textbf{Structure of design rules.}
A design rule follows an IF/THEN form. The IF clause defines a
\textit{quantifier} that selects relevant program elements, while the
THEN clause defines a \textit{constraint} those elements must satisfy
~\cite{mehrpour2019active}. A checker reports an element when it
matches the quantifier but not the constraint, turning a design
decision into continuously executable documentation.

\textbf{Concrete CrowdCode rule.}
Rules range from statement-level conditions to class-level and
cross-file architectural constraints. Listing~\ref{lst:dto-rule}
shows an actual instance of CrowdCode's DTO-constructor rule. The rule
applies to classes extending \texttt{DTO} and requires them to declare
a zero-argument constructor so that Jackson can instantiate them during
deserialization.

In the evaluated CrowdCode snapshot, \texttt{AnswerDTO} violates the
rule because its only explicit constructor requires two arguments.
\texttt{WorkerInFirebase} provides a same-project satisfying example
because it declares an empty zero-argument constructor.

\begin{lstlisting}[
  language=Java,
  caption={Abridged CrowdCode excerpts showing a violation and a satisfying example of the DTO-constructor rule.},
  label={lst:dto-rule}
]
// src/com/crowdcoding/dto/ajax/questions/AnswerDTO.java
public class AnswerDTO extends DTO {
  public long questionId;
  public String text;

  public AnswerDTO(long questionId, String text) {
    this.questionId = questionId;
    this.text = text;
  }
}

// src/com/crowdcoding/dto/firebase/WorkerInFirebase.java
public class WorkerInFirebase extends DTO {
  // fields omitted

  public WorkerInFirebase() { }

  // parameterized constructor omitted
}
\end{lstlisting}

\section{Related Work}
\label{sec:related-work}
% Add section text here.

\subsection{Static Analysis for Defect Detection and Repair}
\label{sec:related-static-analysis}
% Add subsection text here.
Static analysis is central to software quality, but many tools provide
limited guidance for resolving reported issues
~\cite{nachtigall2022usability}. Sorald repairs selected SonarQube
violations through rule-specific AST transformations
~\cite{etemadi2022sorald}, while CORE uses a proposer-ranker LLM
architecture~\cite{wadhwa2024core}. ActiveDocumentation represents
project decisions as executable srcML/XPath rules
~\cite{mehrpour2019active}, and RulePad supports their authoring
~\cite{mehrpour2020rulepad}. DesignFix advances this checker-centered
foundation with project-grounded explanations, generated repairs, and
developer-controlled review, without requiring a hand-written
transformation for every supported rule.

\subsection{LLM-Based Automated Program Repair}
\label{sec:related-apr}
% Add subsection text here.
LLM-based program repair spans zero-shot infilling
~\cite{xia2022alpharepair}, conversational use of test feedback
~\cite{xia2024chatrepair}, and autonomous repair planning
~\cite{bouzenia2025repairagent}. Other systems improve repair through
explicit decomposition or retrieved examples
~\cite{joshi2023ring,wang2023rapgen,jin2023inferfix}. SRepair uses one
model to produce natural-language repair guidance and another to
generate a patch~\cite{xiang2024srepair}. DesignFix targets a distinct
repair signal: a checker-confirmed, project-specific design rule. It
grounds generation in same-project satisfying examples and separates
repair reasoning from structured extraction through two calls to the
same model.

\subsection{LLM-Assisted Developer Tools and Human-AI Interaction}
\label{sec:related-hai}
% Add subsection text here.
AI coding assistants can accelerate implementation, but developers
incur substantial costs when evaluating generated code
~\cite{liang2024survey,barke2023grounded}. Incorrect suggestions can
also produce misplaced confidence~\cite{perry2023insecure}, making
workflow compatibility and review support important for adoption
~\cite{russo2024navigating}. Diff-based interfaces are particularly
effective for reviewing multi-line changes
~\cite{vaithilingam2023intellicode}. DesignFix incorporates these
findings through visible explanations, explicit patch review, and
developer-controlled application.

\subsection{Prompting Strategies for Code Tasks}
\label{sec:related-prompting}
% Add subsection text here.
Structured and plan-first prompting improve code generation by
separating reasoning from implementation
~\cite{li2024scot,jiang2024selfplanning}, while explanation and
self-refinement can help models identify defects in generated code
~\cite{chen2024selfdebug,madaan2023selfrefine}. Retrieved documentation
and demonstrations further improve performance on unfamiliar APIs and
project conventions~\cite{zhou2023docprompting,nashid2023cedar}.
DesignFix combines these ideas by grounding repair in a same-project
satisfying example and separating repair reasoning from structured
output generation.

\section{DesignFix: System Design}
\label{sec:system-design}
% Add section text here.
DesignFix transforms a checker-reported violation into a
project-grounded, developer-reviewable repair through three stages:
context extraction, two-step LLM generation, and interactive review.

\subsection{System Architecture}
\label{sec:system-architecture}
% Add subsection text here.

\begin{figure*}[t]
  \centering
  \includegraphics[width=\textwidth]{DesignFix_System_Diagram_v5.png}
  \caption{DesignFix pipeline.
  \protect\circmark{A} Design rules expressed as XPath queries detect
  violations in the XML AST, producing sets of rule violations and
  examples of the rule being followed.
  \protect\circmark{B} Surrounding context is extracted by retaining
  the violation's enclosing method body and stripping sibling methods
  to signatures.
  \protect\circmark{C} Context elements populate a prompt template
  processed through a two-step LLM pipeline separating free-form
  reasoning from structured JSON extraction.
  \protect\circmark{D} The developer reviews the suggested fix and
  explanation via a side-by-side diff view before accepting or rejecting.}
  \Description{A four-stage workflow showing design-rule detection,
  extraction of surrounding code context, a two-step language-model
  pipeline, and developer review and application of the proposed fix.}
  \label{fig:architecture}
\end{figure*}

Figure~\ref{fig:architecture} illustrates the DesignFix architecture.
When a developer edits code, DesignFix uses srcML
~\cite{collard2013srcml} to represent Java source as XML and checks each
design rule against that representation. Each rule is stored in
\texttt{ruleTable.json} as a natural-language description, a source-file
scope, and quantifier and constraint XPath queries.

Elements matching the quantifier but not the constraint are reported as
violations. Elements matching both queries are satisfying examples.
Listing~\ref{lst:rule-json} reproduces the relevant fields of the actual
CrowdCode DTO-constructor rule. Its quantifier selects classes extending
\texttt{DTO}; its constraint selects those classes that also declare a
constructor with zero parameters. Consequently, \texttt{AnswerDTO}
matches the quantifier but not the constraint, while
\texttt{WorkerInFirebase} matches both.

The stored natural-language description additionally states that the
zero-argument constructor should have an empty body. The executable
XPath checks the presence of a zero-argument constructor but does not
inspect whether its body is empty.

For each reported violation, DesignFix supplies one satisfying match as
the project-specific example. The implementation first uses a satisfying
match in the currently open file when one exists; otherwise, it uses the
first satisfying match in the checker's loaded-file order. It performs
no semantic ranking among candidates. For the run illustrated here, the
selected example was \texttt{WorkerInFirebase.java}.

DesignFix builds on the established rule representation and checking
mechanisms of ActiveDocumentation~\cite{mehrpour2019active} and
RulePad~\cite{mehrpour2020rulepad}. Its contribution begins at the
checker-to-repair boundary, where detected violations are transformed
into project-grounded explanations, candidate repairs, and a
developer-review workflow.

\begin{lstlisting}[
  caption={Relevant fields of the CrowdCode DTO-constructor rule in \texttt{ruleTable.json}.},
  label={lst:rule-json},
  basicstyle=\ttfamily\scriptsize,
  breaklines=true
]
{
  "index": "9",
  "title":
    "DTOs must have a zero-argument constructor",
  "description":
    "IF a class is a DTO object \nTHEN it needs
     to have a zero-argument constructor with
     an empty body.",
  "checkForFilesFoldersConstraints": "INCLUDE",
  "checkForFilesFolders": [
    "src/com/crowdcoding/dto"
  ],
  "processFilesFolders": "WITHIN",
  "quantifierXPathQuery": [
    "//src:unit/src:class[
       src:super/src:extends/src:name/text()=\"DTO\"]"
  ],
  "constraintXPathQuery": [
    "//src:unit/src:class[
       src:super/src:extends/src:name/text()=\"DTO\"
       and src:block/src:constructor[
         count(src:parameter_list/src:parameter)=0]]"
  ]
}
\end{lstlisting}
Upon detecting a violation, the context extraction module assembles the relevant information — the rule description, a rule-satisfying example from the same codebase, the violating code, and the surrounding file context — and passes it to the two-step LLM pipeline described below. The resulting fix and explanation are then presented to the developer through an interactive diff-based review interface.

\subsection{Context Extraction for LLM Prompting}
\label{sec:context-extraction}
A central challenge in leveraging LLMs for design rule repair is 
determining what prompt information to include in the prompt. Through iterative 
experimentation (Section~\ref{sec:evaluation-one}), we identified 
five types of context that improve LLM performance:

\textbf{Design Rule Description.} The natural language description of 
the design rule, explaining what constraint is imposed and when it 
applies. This provides the LLM with semantic understanding of the 
rule's intent.

\textbf{Rule-Satisfying Example.} A code snippet from the same codebase 
that correctly follows the rule, along with its file path, demonstrating 
the expected pattern. We found that providing a project-specific example, 
rather than a generic one, substantially improved fix quality by exposing 
local conventions and naming patterns.

\textbf{Violation Snippet.} The specific lines of code (typically 1--2 
lines) that trigger the design rule violation, as identified by the 
XPath query evaluation. %This snippet pinpoints what is wrong without 
including extraneous code. 
The program element requiring modification may differ from the reported violation site.

\textbf{Surrounding Context.} The enclosing class and method structure 
around the violation, extracted from the XML AST by traversing upward 
from the violation node. The extraction strategy varies by the type of 
code element violated:

\begin{itemize}
\item \textit{Statements within a method:} The full body of the 
enclosing method is retained, while sibling methods in the same class 
are reduced to their signatures (bodies removed). Field declarations 
and the class declaration are preserved.
\item \textit{Class-level elements} (e.g., fields, annotations): The 
enclosing class is returned with all method bodies stripped to 
signatures only, preserving the structure of the class.
\item \textit{Static utility classes:} The entire class is included, 
as static classes tend to be compact with interdependent methods.
\end{itemize}

\noindent This strategy provides the LLM with sufficient structural 
context---the class shape, field declarations, and method 
signatures---while filtering out irrelevant implementation details from 
sibling methods. %Listing~\ref{lst:context_extraction} in the Appendix illustrates this 
%extraction for a statement-level violation.
\begin{comment}
\begin{lstlisting}[language=Java,caption={Surrounding context 
extraction example. For a violation inside \texttt{execute()}, the 
extracted context retains the full method body and class structure 
but strips sibling method bodies to signatures 
only.},label={lst:context_extraction}]
// Full class (original)
public class CreateFunction extends Command {
  private String name;
  private String description;
  
  public void execute() {
    // ... 15 lines (contains violation)
  }
  
  public void validate() {
    // ... 10 lines (unrelated)
  }
  
  public static CreateFunction create() {
    // ... 5 lines (unrelated)
  }
}

// Extracted surrounding context
public class CreateFunction extends Command {
  private String name;
  private String description;
  
  public void execute() {
    // ... 15 lines (full body retained)
  }
  
  public void validate() { ... }
  
  public static CreateFunction create() { ... }
}
\end{lstlisting}    
\end{comment}


\textbf{File Context.} The complete file containing the reported violation, enabling the LLM to understand the local context in which the rule was violated, including package declarations, imports, and class definitions. For cross-file rules, the repair site may reside in another file.

Although the file context contains the complete violation file, the
surrounding context is not redundant with it: it is a filtered
structural view in which package and import declarations are removed
and sibling method bodies are reduced to signatures, giving the model
the class shape relevant to the violation without distracting detail.
Presenting this scoped view as a separately labeled section, rather
than relying on markers embedded within the full file, signals which
region the model must understand but not modify; this explicit
separation measurably reduced unintended edits to context code
(Finding~3, Section~\ref{sec:evaluation-one-findings}). The full file
complements it by supplying the integration details the skeleton
omits, such as imports and package placement.

% Add subsection text here.

\subsection{Two-Step Co-Generative Prompting}
\label{sec:two-step-prompting}
% Add subsection text here.
Our evaluation (Section~\ref{sec:evaluation-one}) found that a two-step prompting approach, which separates free-form reasoning from structured output generation, outperforms single-shot prompts.

\textbf{Step 1: Unstructured Reasoning.} The first prompt contains the five inputs described above and instructs the model to produce a free-form response explaining the violation and a fix. The prompt explicitly encourages the model to ``take as much time as needed'' to reason through the problem without worrying about output format. This step uses a temperature of 0.50, which our evaluation (Section~\ref{sec:evaluation-one}) identified as optimal for balancing reasoning quality and consistency.

Figure~\ref{fig:architecture} shows the structure of the Step~1 prompt. Each context element is enclosed in explicitly labeled delimiters, enabling the model to distinguish the reported violation file, the structural context surrounding the violation, and the specific code that triggered the rule. When required by the rule, the generated suggestion may identify a repair location in another file.
\begin{comment}
\begin{lstlisting}[caption={Step 1 prompt template with labeled 
context sections.},label={lst:prompt_template},
basicstyle=\ttfamily\scriptsize]
You are assisting with enforcing the following 
design rule:
{rule description}

Here is a code example that follows the rule:
{rule-satisfying example}
Example file path: {example file path}

<<<ORIGINAL_FILE_PATH>>>
{violation file path}
<<<END_ORIGINAL_FILE_PATH>>>

<<<ORIGINAL_FILE_CONTENT>>>
{full file content}
<<<END_ORIGINAL_FILE_CONTENT>>>

Here is the surrounding class/method structure 
around the violation, showing the enclosing context:
<<<SURROUNDING_CONTEXT>>>
{extracted surrounding context}
<<<END_SURROUNDING_CONTEXT>>>

<<<VIOLATION_SNIPPET>>>
{violation snippet}
<<<END_VIOLATION_SNIPPET>>>

Rewrite the original file so it satisfies the rule 
while preserving every unrelated line verbatim.
\end{lstlisting}    
\end{comment}
The explicit separation between
\texttt{SURROUNDING\_CONTEXT} and \texttt{VIOLATION\_SNIPPET} is a
deliberate design choice: the labeled delimiters signal that the
surrounding context is provided for understanding, not for
modification, which measurably reduced unintended edits to context
code (Finding~3, Section~\ref{sec:evaluation-one-findings}).


\textbf{Step 2: Structured Extraction.} The second prompt feeds back both the original prompt and the unstructured response, asking the model to restructure the answer into JSON with three fields: \texttt{\{explanation, code, fileName\}}. This step uses a lower temperature (0.2) to ensure precise, parseable output.
This separation allows the model to reason freely without formatting constraints, provides implicit self-review during restructuring, and ensures reliable output.

\subsection{Interactive Suggestion and Acceptance Workflow}
\label{sec:interactive-workflow}
% Add subsection text here.
Figures~\ref{fig:ui-a}--\ref{fig:ui-b} illustrate this workflow using
a concrete example from the CrowdCode codebase, where a DTO class
violates the running-example rule of Listing~\ref{lst:dto-rule}: it is
missing the required zero-argument constructor.
Upon receiving the structured JSON response, the system displays the suggestion in a dedicated UI panel showing the \textbf{Suggested Fix Location} (file path), the \textbf{Suggested Fix} (code), and the \textbf{Explanation} of changes. The developer reviews this information and may click \textbf{Accept Fix} to proceed. Clicking Accept Fix triggers a side-by-side diff view for the suggested repair file, which may differ from the file containing the reported violation. The left tab displays the current contents of the suggested repair file, and the right tab shows the proposed modification highlighted within its surrounding code. Accept Change applies the modification to the suggested repair file, while Reject Change dismisses the suggestion.
\begin{figure}[t]
  \centering
  \includegraphics[width=\columnwidth]{1_anno.png}
  \caption{A developer encounters a design rule violation.
  The rule checker surfaces
  {\ding{192}}~the rule description explaining the IF/THEN constraint,
  {\ding{193}}~categorization tags for the violation type,
  {\ding{194}}~the specific code snippets that violate the rule, and
  {\ding{195}}~a Fix button that invokes the DesignFix LLM pipeline.}
  \Description{A DesignFix interface displaying a design-rule
  description, rule-category tags, violating code snippets, and a button
  for requesting an LLM-generated repair.}
  \label{fig:ui-a}
\end{figure}

\begin{figure}[t]
  \centering
  \includegraphics[width=\columnwidth]{2_anno_alt.png}
  \caption{After clicking Fix, DesignFix returns a structured suggestion
  showing
  {\ding{192}}~the target file path,
  {\ding{193}}~the suggested code fix with additions marked by
  \texttt{+},
  {\ding{194}}~a natural-language explanation grounded in
  project-specific context, and
  {\ding{195}}~an Accept Fix button to proceed to review.}
  \Description{A DesignFix suggestion panel displaying the target file,
  proposed patch, explanation, and a control for proceeding to review.}
  \label{fig:ui-b}
\end{figure}

\section{Evaluation I: Prompting Strategies}
\label{sec:evaluation-one}
% Add section text here.
To understand how effectively LLMs can assist with design rule violation repair, we systematically investigated different prompting strategies and measured their effectiveness across multiple metrics. This evaluation addresses RQ1 and RQ2.

\subsection{Methodology}
\label{sec:evaluation-one-method}
% Add subsection text here.
\textbf{Prompt Development.} We developed prompts iteratively, starting with a minimal baseline and progressively adding context and constraints, as discussed in Table~\ref{tab:prompts}. We developed prompts using a subset of three design rule violations from one codebase (CrowdCode), then evaluated the finalized prompts on the full dataset. These three development violations were included in the 45-violation evaluation dataset to maximize data usage. This development/evaluation split ensured that our prompt refinements generalized beyond the specific examples used during development.

\begin{table*}[t]
  \caption{Prompt variants evaluated in the LLM proficiency study.
  Each variant builds incrementally on the previous variant.}
  \label{tab:prompts}
  \centering
  \begin{tabular}{lll}
    \toprule
    \textbf{Prompt} & \textbf{Description} & \textbf{Tests} \\
    \midrule
    Baseline & Rule description + violating code only & Minimal context \\
    Prompt 1 & + Output only fix lines (no prose) & Succinct patches \\
    Prompt 2 & + Rule-satisfying example from codebase & Positive examples \\
    Prompt 3 & + JSON output \{explanation, code\} & Structured output \\
    Prompt 4 & + Separately labeled surrounding class/method structure
             & Context integration \\
    Prompt 5 & + Include surrounding code in output
             & Context preservation \\
    \midrule
    Final & Two-step: free-form reasoning $\rightarrow$ JSON
          & Reasoning/format separation \\
    \bottomrule
  \end{tabular}
\end{table*}

\textbf{Dataset.}
We evaluated 15 design rules from two rule corpora:

\begin{itemize}
\item \textit{CrowdCode:} Eight project-specific rules spanning command
patterns, entity persistence, and artifact relationships. The rules
originated with CrowdCode developers and were subsequently documented
and encoded for automated checking in prior work
~\cite{mehrpour2019active}. CrowdCode is an open-source Java
microtask-programming system with approximately 9,000 lines across 107
classes. We seeded violations into compiling project snapshots.

\item \textit{Android Architectural Rules:} Seven rules derived from
Android guidance, covering DTO, repository, and persistence patterns.
\end{itemize}

For each rule, we evaluated three violation instances, yielding 45
cases spanning missing annotations, incorrect signatures, and invalid
class relationships.
\begin{comment}
\textbf{Dataset.} We gathered design rules and corresponding violations from two distinct codebases to ensure generalizability:

\begin{itemize}
\item \textit{CrowdCode:} an open-source cloud IDE for microtask
programming, in which crowds build software by completing short,
self-contained microtasks coordinated through a graph of artifacts
such as functions and tests~\cite{latoza2014microtask,
latoza2019microtaskTSE}. Its back-end comprises approximately 9,000
lines of Java across 107 classes, implemented on Google App Engine
with the Objectify persistence framework, and is publicly
available.\footnote{\url{https://github.com/devuxd/CrowdCode}; the
Java implementation used in this work resides on the repository's
earlier branches.} We used 8 design rules covering command patterns,
entity persistence, and artifact relationships, drawn verbatim from
the 15 rules extracted from this codebase in prior
work~\cite{mehrpour2019active}. Violations were introduced by
editing CrowdCode source files to recreate realistic developer
mistakes, yielding project versions that compile while containing
the seeded violation.
\item \textit{Android Architectural Rules:} A set of 7 rules derived from Android best practices covering data transfer objects (DTOs), repository patterns, and persistence patterns. These reflect modern Android design recommendations.
\end{itemize}

For each of the 15 design rules, we obtained or constructed 3 instances of violating code, resulting in 45 violation examples. Violations were selected to represent realistic developer mistakes, including missing annotations, incorrect method signatures, and improper class relationships.
\end{comment}

\textbf{LLM Configurations.}
We evaluated the seven variants with GPT-4o snapshots from August 6
(\textit{gpt-4o-aug}) and November 20, 2024
(\textit{gpt-4o-nov}). Each prompt, model, and generation-temperature
combination (0.5, 0.75, and 1.0) was sampled five times. Final used the
tested temperature for reasoning and 0.2 for structured extraction.
Across 45 violations, this design produced 9,450 end-to-end repair
trials.

\begin{comment}
\textbf{LLM Configurations.} We used OpenAI's gpt-4o models, comparing two versions: gpt-4o-2024-11-20 and gpt-4o-2024-08-06. Comparing versions helps us understand whether prompt effectiveness varies across model updates. To account for model nondeterminism, we ran each prompt with three temperature settings (0.5, 0.75, 1.0) and generated 5 outputs per configuration. For the final prompt, while our evaluation tested temperatures 0.5, 0.75, and 1.0 for the reasoning step to determine optimal creativity/precision balance, we use a lower temperature (0.2) for the structuring step in production to maximize JSON formatting reliability—a parameter choice validated through preliminary pilot testing. This resulted in up to 15 attempts per scenario, enabling statistical analysis of reliability. These temperature values were selected because they span the commonly recommended range for code generation tasks. In total, our experiment involved:

45 violations $\times$ 7 prompt types $\times$ 2 model versions $\times$ 3 temperatures $\times$ 5 trials = 9,450 LLM runs. For the two model versions, we evaluated two gpt-4o models.
Throughout this paper, we will refer 
\begin{itemize}
    \item gpt-4o-2024-08-06 as gpt-4o-aug  and
    \item gpt-4o-2024-11-20 as gpt-4o-nov
\end{itemize} 
\end{comment}

\textbf{Evaluation Metrics.} We assessed LLM suggestions along three dimensions:

\textit{Code Fix Success Rate:} Whether the suggested patch resolves
the reported design-rule violation. We applied each patch and reran the
rule checker; a patch passed when the checker no longer reported the
violation.

\textit{Surrounding Code Pass Rate:} Whether the proposed patch avoids
changes that are unrelated or unnecessary to repairing the reported
violation. Necessary supporting edits, such as adding an import or
modifying a separate repair file, are permitted. This metric captures
preservation of unrelated code and is assessed independently of whether
the patch resolves the violation.

\textit{Explanation Correctness:} Whether the explanation accurately
identifies the violated rule, correctly explains why the reported code
violates it, and remains consistent with the proposed patch.

\textbf{Labeling Procedure.}
\textit{Code Fix Success} was determined by rerunning the rule checker after
applying each patch. Explanation Correctness and Surrounding Code Pass
Rate were manually labeled by the first author using rubrics fixed
before labeling began.

\textit{Explanation Correctness} passed only if the explanation
(a) identified the rule that was violated, (b) correctly explained why
the code violated that rule, and (c) was consistent with the proposed
patch.

\textit{Surrounding Code Pass Rate} was assessed independently of Code
Fix Success. It passed when a diff against the original project snapshot
contained no unrelated or unnecessary modifications. Changes required
to implement the attempted repair---including necessary imports,
declarations, or edits in a separate repair file---were permitted. A
patch could therefore preserve surrounding code even if it failed to
resolve the reported violation.


\subsection{Results}
\label{sec:evaluation-one-results}
% Add subsection text here.


\begin{figure}[t]
  \centering
  \includegraphics[width=\columnwidth]{codefix_passrate.png}
  \caption{Code Fix Pass Rate across five prompt variations and three
  temperature parameters. Results from the Baseline prompt and Prompt~1
  are omitted for brevity. Labels are abbreviated as P2--P5 for
  Prompts~2--5 and F for the Final prompt, with \emph{nov} and
  \emph{aug} denoting the GPT-4o November and August model variants.}
  \Description{A grouped bar chart comparing code-fix pass rates across
  prompt variants, model versions, and temperature settings.}
  \label{fig:code-fix}
\end{figure}

\textbf{Code Fix Success Rate.} Figure~\ref{fig:code-fix} presents code fix success rates across prompt variants, temperatures, and model versions. The Baseline prompt represents unassisted foundation-model use: it supplies only the rule description and the violating snippet, approximating what a developer obtains by pasting a checker warning and the flagged code into a conversational LLM such as ChatGPT, without any of DesignFix's context augmentation. This baseline achieved only 62--64\% success at temperature 0.5.
Adding code-only output constraints (Prompt 1) improved success to 68-70\%, and adding rule-following examples (Prompt 2) further improved success to 76-77\%. The Final Prompt (two-step co-generative) achieved the highest success rates across all conditions. At temperature 0.5, the Final Prompt achieved 88\% success with gpt-4o-nov and 90\% with gpt-4o-aug. At temperature 0.5, this represents an exact
26-percentage-point improvement over the matched Baseline for both
model snapshots and an 8--10-point improvement over Prompt~5, its
closest single-step counterpart. This matched 26-point gap quantifies
the combined difference associated with DesignFix's context
augmentation and two-step generation--- same-codebase rule-satisfying examples, scoped surrounding context,
and two-step generation --- over handing the same violation directly to a foundation model.


\begin{comment}
Lower temperatures generally produced better results, with success rates declining by 2-5 percentage points as temperature increased from 0.5 to 1.0. This aligns with expectations that lower temperatures produce more deterministic, focused outputs suitable for code generation tasks.
\end{comment}

\begin{comment}
\begin{table}[t]
\caption{Explanation Pass Rate (\%)}
\label{tab:explanation}
\centering
\begin{tabular}{llccc}
\toprule
\textbf{Prompt} & \textbf{Model} & \textbf{0.5} & \textbf{0.75} & \textbf{1.0} \\
\midrule
Prompt 3 & gpt-4o-nov & 76 & 76 & 73 \\
Prompt 3 & gpt-4o-aug & 77 & 77 & 72 \\
Prompt 4 & gpt-4o-nov & 75 & 74 & 70 \\
Prompt 4 & gpt-4o-aug & 76 & 75 & 71 \\
Prompt 5 & gpt-4o-nov & 82 & 74 & 80 \\
Prompt 5 & gpt-4o-aug & 84 & 76 & 76 \\
\midrule
Final & gpt-4o-nov & \textbf{89} & 87 & 85 \\
Final & gpt-4o-aug & 88 & 87 & 82 \\
\bottomrule
\end{tabular}
\end{table}
\end{comment}

\begin{figure}[t]
  \centering
  \includegraphics[width=\columnwidth]{explanation_passrate.png}
  \caption{Explanation Pass Rate across four prompt variations and
  three temperature parameters. Prompts earlier than Prompt~3 are
  omitted because they did not generate explanations.}
  \Description{A grouped bar chart comparing explanation pass rates
  across prompt variants, model versions, and temperature settings.}
  \label{fig:explanation-passrate}
\end{figure}

\textbf{Explanation Pass Rate.} Figure~\ref{fig:explanation-passrate} presents explanation pass rates for prompts that generated explanations (Prompts 3-5 and Final). The Final Prompt achieved 89\% explanation correctness with gpt-4o-nov at temperature 0.5, compared to 76-77\% for Prompt 3.
Observed explanation failures included misidentifying the applicable
rule, incorrectly explaining why the code violated it, and producing
an explanation inconsistent with the proposed patch. The two-step approach reduced all three error types, with the largest reduction in explanation-fix mismatches.

\begin{comment}
\begin{table}[t]
\caption{Surrounding Code Pass Rate (\%)}
\label{tab:surrounding}
\centering
\begin{tabular}{llccc}
\toprule
\textbf{Prompt} & \textbf{Model} & \textbf{0.5} & \textbf{0.75} & \textbf{1.0} \\
\midrule
Prompt 4 & gpt-4o-nov & 76 & 72 & 70 \\
Prompt 4 & gpt-4o-aug & 86 & 83 & 72 \\
Prompt 5 & gpt-4o-nov & 80 & 75 & 80 \\
Prompt 5 & gpt-4o-aug & 78 & 81 & 79 \\
\midrule
Final & gpt-4o-nov & \textbf{90} & 88 & 86 \\
Final & gpt-4o-aug & 89 & 85 & 85 \\
\bottomrule
\end{tabular}
\end{table}
\end{comment}

\begin{figure}[t]
  \centering
  \includegraphics[width=\columnwidth]{surrounding_code_passrate.png}
  \caption{Surrounding Code Pass Rate across three prompt types,
  three temperature settings, and two model versions.}
  \Description{A grouped bar chart comparing surrounding-code
  preservation rates across prompt variants, model versions, and
  temperature settings.}
  \label{fig:surrounding-passrate}
\end{figure}

\textbf{Surrounding Code Pass Rate.} Figure~\ref{fig:surrounding-passrate} presents Surrounding Code Pass
Rates from Prompt~4 onward, when preservation became an explicit
evaluation concern. Prompt~4 supplied surrounding context as input,
whereas Prompt~5 and Final generated integrated repair code. The final prompt achieved 90\% at temperature 0.5 with gpt-4o-nov, substantially outperforming Prompts 4 (76\%) and 5 (80\%).
This metric proved important because earlier prompts sometimes made unintended modifications to surrounding code when asked to include it in the output, such as removing necessary import statements. Finding~3 quantifies these unintended-modification categories and the reduction achieved by the two-step approach.


\subsection{Key Findings}
\label{sec:evaluation-one-findings}
% Add subsection text here.
\textbf{Finding 1: Rule-following examples substantially improve fix
quality.} Across the six matched model--temperature conditions,
Prompt~2, which included a rule-satisfying example, exceeded Prompt~1
by 7--10 percentage points for Code Fix Success. At temperature 0.5,
the gains were 8 points for gpt-4o-nov and 7 points for gpt-4o-aug.
The example provides a concrete implementation of the rule that exposes
project-specific naming, import, framework, and structural conventions.

\textbf{Finding 2: Structured output requires careful prompting.} Requesting JSON output (Prompt 3) sometimes led to format deviations or extra text outside the JSON structure in 15\% of cases. Common issues included markdown formatting around JSON blocks, explanatory preambles before the JSON, and incomplete JSON due to token limits. The two-step approach mitigated this by separating reasoning from formatting, reducing format errors to under 3\%. This suggests that constraining output format early in the generation process interferes with the model's reasoning capabilities.

\textbf{Finding 3: Surrounding context improves fix integration when 
explicitly separated from the violation.} Providing surrounding code 
as input (Prompt~4) helped the model understand structural context and 
produce fixes that integrated correctly, improving integration success 
by 3--14 percentage points across model versions. When 
surrounding context was not clearly delineated from the violation 
code, the model sometimes made unintended modifications: renaming 
variables to ``improve'' code style (in 12\% of cases), adding 
unsolicited comments (8\%), or refactoring unrelated code (4\%). 
Introducing explicitly labeled delimiters 
(\textit{<<<SURROUNDING\_CONTEXT>>>} and 
\textit{<<<VIOLATION\_SNIPPET>>>}) reduced unintended modification 
rates from 24\% to 10\% by signaling that surrounding code is 
provided for comprehension, not for modification. The two-step 
approach further reduced this rate by encouraging the model to 
extract only the relevant fix from its initial reasoning.


%\textbf{Finding 3: Surrounding context is a double-edged sword.} Providing surrounding code as input (Prompt 4) helped the model understand context and produce fixes that integrated correctly, improving integration success by 3-14 percentage points across model versions, from Prompt 4 to the final prompt. However, requesting surrounding code in the output sometimes led to unintended modifications: the model would ``improve'' variable names, add unsolicited comments, or refactor unrelated code. The two-step approach reduced unintended modification rates from 24\% to 10\% by encouraging the model to extract only the relevant fix from its initial reasoning.

\textbf{Finding 4: The two-step approach consistently outperforms
single-shot prompts.} Final achieved the highest observed rate in all
18 model--temperature--metric combinations. At temperature 0.5, it
improved over Prompt~5 by 8--10 points for Code Fix Success, 4--7
points for Explanation Correctness, and 10--11 points for Surrounding
Code Pass Rate. The trade-off is a second API call, which added a few
seconds in the observed runs.

\textbf{Finding 5: Lower temperatures generally produce stronger
results.} Temperature 0.5 produced the highest or tied-highest pass
rate in 27 of the 28 prompt--model--metric series. The exception was
Prompt~5's Surrounding Code Pass Rate with gpt-4o-aug, which reached
81\% at temperature 0.75 compared with 78\% at 0.5. We therefore use
0.5 as the best-performing tested reasoning temperature.

\section{Evaluation II: User Study}
\label{sec:user-study}
% Add section text here.
To understand how DesignFix affects developer task performance, we conducted a within-subjects user study comparing two workflows: a without DesignFix condition and a DesignFix condition. This evaluation addresses RQ3.

\subsection{Methodology}
\label{sec:user-study-method}
% Add subsection text here.
\textbf{Participants.} We recruited 14 participants with diverse software engineering backgrounds through academic and professional networks. The cohort included PhD students in Computer Science (n=10), professional software engineers (n=3), and a recent CS graduate (n=1). Their specializations spanned software development, machine learning and NLP, network security, VR development, FinTech, and university lecturing. Professional industry experience ranged from 0 to 12 years, with a median of 2.0 years. All participants had experience with Java and were familiar with object-oriented design patterns, though none had prior experience with the specific codebase used in the study. This selection ensured participants had sufficient technical skill to complete the tasks while also being unfamiliar with the codebase, simulating the common scenario of joining a new project.

\textbf{Tasks and task selection.}
Participants worked with the CrowdCode codebase described in
Section~\ref{sec:evaluation-one-method}. The candidate pool consisted
of the 15 project-specific CrowdCode rules previously elicited and
encoded in the ActiveDocumentation work
~\cite{mehrpour2019active}. We used these existing rule definitions
without modification.

Both authors purposively selected eight study rules after inspecting
the Evaluation~I outcomes. We included rules associated with
unsuccessful repairs and then filled the remaining study set with the
most difficult rules associated with successful repairs. Difficulty was
assessed using the expected number of lines requiring modification, the
amount of codebase understanding needed beyond the reported violation,
and the architectural or cross-file reasoning required. The authors
then organized the eight rules into four pairs judged comparable using
these criteria.

Each task required the participant to repair one assigned seeded
violation while preserving the surrounding program and CrowdCode's
project-specific conventions. This was an outcome-informed,
purposively selected task set rather than a random sample of CrowdCode
rules.

\textbf{Study Design.}
We used a within-subjects design in which each participant completed
two design-rule-repair tasks: one in the checker-centered condition
(without DesignFix) and one in the DesignFix condition.

In both conditions, participants used Visual Studio Code, the same
CrowdCode repository, and the ActiveDocumentation checker. The checker
displayed the rule description, categorization tags, violating snippet,
and file path. Clicking the reported violation opened the corresponding
file in Visual Studio Code at the reported location. Thus, both
conditions received violation-site localization, although the reported
location was not necessarily the program element that needed to be
modified.

In the checker-centered condition, participants used the checker
output, repository, and permitted browser-based resources to determine
the repair location and implementation. They did not receive
DesignFix's same-project satisfying example, extracted repair context,
generated explanation, proposed patch, or review interface.

In the DesignFix condition, participants received the same repository,
rule material, checker output, and file-opening affordance together
with the complete DesignFix bundle: a same-project satisfying example,
extracted context, generated explanation, repair-site guidance,
proposed patch, and review interface. Neither condition was given the
correct repair for a study task.

Browser-based conversational LLMs, including ChatGPT, Claude.ai, and
Gemini, were permitted in the checker-centered condition. In-IDE coding
assistants were prohibited in both conditions. Their uses included requesting definitions and conceptual
clarifications, debugging or validating tentative changes, and, in
some cases, generating a candidate repair.

The checker-centered comparator represents the existing workflow that
DesignFix augments; it is not a no-support baseline. The comparison
therefore estimates the effect of the integrated DesignFix bundle and
does not isolate the contribution of any single component or compare
LLM access with no LLM access.

\textbf{Procedure.}
The study was conducted remotely through video conferencing and screen
sharing. Before either task, every participant received the same
10-minute introduction to design rules and CrowdCode, followed by the
same 5-minute tutorial covering both the checker-centered and DesignFix
workflows. The instruction did not disclose the correct repairs for the
study tasks.

During the DesignFix task, the reasoning step used
\textit{gpt-4o-2024-11-20} at temperature 0.5, and the structured
extraction step used temperature 0.2.

Each task had a 20-minute limit. Before every task, the facilitator
restored the working repository from the same seeded backup using
Robocopy. The backup contained all seeded study violations, and
participants were instructed to address only their assigned violation.
After each task, the working copy was discarded and restored before
subsequent use. This prevented edits made by a preceding participant or
task from carrying forward.

Participants were instructed to repair the assigned violation while
preserving the project's existing conventions. After each task, they
completed the NASA Task Load Index (NASA-TLX)
~\cite{hart1988development} and answered semi-structured questions
about their strategies, experience, and difficulties. Sessions were
screen-recorded with participant consent.

\begin{comment}
\textbf{Violation and Repair Locations}
The design-rule checker localized violating code. However, the violating snippet was not necessarily the location that needed to be modified. Some rules captured cross-file architectural relationships in which a violation was observed in one file but could only be resolved by changing another file. We therefore distinguish violation localization, which identifies the code that fails the rule, from repair-site localization, which identifies the file or program element that must be modified to restore compliance.

In without DesignFix condition, participants received the violating snippet from the checker but had to determine the appropriate repair location and implementation by exploring the repository. In the DesignFix condition, DesignFix’s explanation and proposed solution could identify both the violating location and the relevant repair location, including cases where these were in different files.
\end{comment}

\textbf{Measures.} We collected:
\begin{itemize}
\item \textit{Task duration:} Time from task start until the participant
declared completion, explicitly stopped working, or reached the
20-minute limit. Because this endpoint includes unsuccessful attempts,
it is not equivalent to time to a correct repair.

\item \textit{Repair correctness:} A submitted repair was labeled
correct only if (a) the checker no longer reported the assigned
violation, (b) the project compiled successfully.
The first author assessed submissions without access to explicit
condition labels, using the checker result, build result, and a diff
against the original seeded snapshot.

\item \textit{NASA-TLX scores:} Six workload dimensions---mental
demand, physical demand, temporal demand, performance, effort, and
frustration---each rated from 1 to 21.
\end{itemize}

\subsection{Quantitative Results}
\label{sec:user-study-quantitative}
% Add subsection text here.

\begin{figure}[t]
  \centering
  \includegraphics[width=\columnwidth]{Task_Completion_Chart_boxplot.png}
  \caption{Task times with and without DesignFix. Boxplots show the
  distributions; points indicate individual participants and lines link
  each participant's two observations.}
  \Description{Paired boxplots and participant-level points comparing
  task-completion times with and without DesignFix.}
  \label{fig:task-completion}
\end{figure}

\textbf{Task Duration.}
Recorded task duration was lower in the DesignFix condition~\ref{tab:time}. Mean
duration was 10.00 minutes (SD=1.80) with DesignFix and 15.71 minutes
(SD=3.02) without DesignFix. A paired-samples t-test found a significant
difference ($t(13)=9.57$, $p<0.001$), with a paired standardized effect
of $d_z=2.56$. The 5.71-minute difference corresponds to a 36.4\%
shorter mean recorded duration. Because the endpoint included
unsuccessful declarations, early stopping, and the 20-minute limit,
these results concern task duration rather than time to a successful
repair.

\begin{table*}[t]
  \caption{Task Duration in Minutes}
  \label{tab:time}
  \centering
  \begin{tabular}{lccccc}
    \toprule
    \textbf{Condition} &
    \textbf{Mean} &
    \textbf{SD} &
    \textbf{Median} &
    \textbf{Min} &
    \textbf{Max} \\
    \midrule
    Without DesignFix & 15.71 & 3.02 & 15.5 & 11 & 20 \\
    With DesignFix    & 10.00 & 1.80 & 10.0 & 7  & 13 \\
    \bottomrule
  \end{tabular}
\end{table*}


\textbf{Repair Correctness.}
Under the correctness criteria described above, 13 of 14
participants (92.9\%) produced a correct repair in the DesignFix
condition, compared with 6 of 14 participants (42.9\%) in the
checker-centered condition. This improvement was statistically significant (exact McNemar's test, $p = 0.016$). Of the seven participants whose outcome changed between conditions, all seven improved (failed without DesignFix but succeeded with it), while none worsened.

%Analysis of the McNemar's contingency table reveals the pattern of change: 7 participants improved (failed without DesignFix but succeeded with it), 6 participants succeeded in both conditions, 0 participants worsened, and 1 participant (P8) failed in both conditions. The absence of any participants who succeeded without DesignFix but failed with it suggests that the tool does not introduce new failure modes for developers who already possess sufficient codebase knowledge.

\begin{table}[t]
  \caption{Error patterns observed in the without DesignFix condition.
  Participants could exhibit multiple patterns.}
  \label{tab:error-patterns}
  \centering
  \small
  \begin{tabular}{lcc}
    \toprule
    \textbf{Error Pattern} & \textbf{Participants} & \textbf{Count} \\
    \midrule
    Syntax/compile errors & P1, P7, P9, P10 & 4 \\
    Wrong method signature & P2, P9, P10 & 3 \\
    Partial fix & P1, P2, P4 & 3 \\
    Misunderstood rule & P3, P4, P7 & 3 \\
    \midrule
    Reached time limit & P4, P9 & 2 \\
    Surrendered early & P8 & 1 \\
    \bottomrule
  \end{tabular}
\end{table}

Table~\ref{tab:error-patterns} summarizes error patterns in the without DesignFix condition; notably, several participants (P4, P9) exhibited multiple error types.
\begin{comment}
The incorrect fixes in the without DesignFix condition exhibited several distinct 
patterns (Table~\ref{tab:error-patterns}). Syntax errors were most common, 
but conceptual errors-wrong method signatures, partial fixes, and 
misunderstanding which architectural layer to edit-were equally prevalent. 
Notably, several participants exhibited multiple error types, with P4 and P9 
each appearing in three categories.
\end{comment}
P3’s experience illustrates the distinction between violation localization and repair-site localization. The checker identified the snippet violating the rule, but compliance required modifying a different file. P3 treated the reported location as the likely repair site and spent approximately 15 of 19 minutes modifying related but incorrect files. In the participant’s DesignFix task, the generated explanation and proposed solution explicitly connected the reported violation to the relevant repair file and program element, enabling the participant to complete the task correctly in 12 minutes.


\begin{figure}[t]
  \centering
  \includegraphics[width=\columnwidth]{NASA_TLX_enlarged.png}
  \caption{Average NASA-TLX ratings across six workload dimensions for
  tasks with and without DesignFix, aggregated over 14 participants.}
  \Description{A chart comparing average NASA-TLX workload ratings
  across six dimensions for the DesignFix and without DesignFix conditions.}
  \label{fig:nasa-tlx}
\end{figure}

\textbf{Cognitive Workload.} NASA-TLX scores were lower across all six
dimensions when using DesignFix (Figure~\ref{fig:nasa-tlx}). To test
significance, we conducted exact Wilcoxon signed-rank tests on each
dimension---appropriate for paired ordinal data with small
samples---and applied Bonferroni correction for six comparisons
(adjusted $\alpha = 0.0083$).
Five of the six dimensions (except Physical Demand which was irrelevant in this context) showed statistically significant reductions
after correction: Mental Demand, Temporal Demand, Performance, Effort, and Frustration, all at ($W=0$, $p<0.001$); all Bonferroni-adjusted $p<0.002$. The $W=0$
indicates that no participant reported higher workload with DesignFix
on any dimension. A Wilcoxon test on the composite (unweighted mean)
score confirmed a significant overall reduction ($10.64$ vs.\ $4.69$;
$W=0$, $p<0.001$).
 
Frustration showed the largest raw reduction, dropping from a mean of
12.07 to 4.43. This was most pronounced among participants who
struggled in the without DesignFix condition. P10, who reported frustration
of 19 without DesignFix, reported only 6 with
assistance. Mental Demand and Effort showed similar large reductions, indicating that DesignFix
substantially reduced the cognitive burden of locating relevant code
and reasoning about unfamiliar patterns. Performance satisfaction also
improved, with sharply decreased standard
deviation, indicating more consistent self-assessed
performance across participants regardless of background.

\subsection{Qualitative Findings}
\label{sec:user-study-qualitative}
% Add subsection text here.
We conducted a detailed analysis of interview transcripts, post-task questionnaire responses, and screen recordings. We identified five major themes.

\textbf{Theme 1: Navigation Dominates Task Difficulty.} The most universally valued feature was DesignFix's ability to identify specific file locations and code segments, mentioned by all 14 participants. Screen recording analysis revealed participants spent an average of 4.9 minutes before their first code modification in the without DesignFix condition, compared to 2.1 minutes with DesignFix.
This reveals that DesignFix's primary value lies in \textit{codebase navigation} rather than code generation. P1 spent the first 10 minutes of the without-DesignFix task searching the web for documentation, finding only unrelated examples instead of guidance specific to this codebase's design rules.
%\textit{``I expected some official spec for this rule, but I kept finding unrelated JPA examples.''}.
With DesignFix, P1 reported 
\textit{``It named the files and basically gave me the patch... No hunting - just edit where it said.''} 


Even expert participants who rejected the automated acceptance workflow benefited substantially. The most experienced participant (P11) manually copied code, yet still achieved 36\% time savings. P11's feedback illuminates this:
\pquote{P11}{with DesignFix}{Even though the tool gave me the patch, I don't trust a button to blindly inject code into the codebase. But it pointed to the exact file and gave a concrete patch, so I didn't have to trace ownership semantics.}

%\textit{``Even though the tool gave me the patch, I don't trust a button to blindly inject code into the codebase. But it pointed to the exact file and gave a concrete patch, so I didn't have to trace ownership semantics.''}



%\pquote{P1}{With DesignFix}{It named the files and basically gave me the patch... No hunting - just edit where it said.}



%P4, who reached the 20-minute time limit without DesignFix, explained:
%\pquote{P4}{Unassisted}{I don't have intuition for this command layer yet. I kept finding classes that looked right but weren't.}

%\textit{``I don't have intuition for this command layer yet. I kept finding classes that looked right but weren't.''} 

%With DesignFix, P4 completed the task in 9 minutes, attributing the improvement to the tool's precise file and line identification.

\textbf{Theme 2: Explanation Quality and Contextual Alignment.} Participants strongly valued explanations that connected the design rule to specific patterns already present in the codebase (mentioned by 13 of 14 participants). P3 reported: 
\pquote{P3}{with DesignFix}{
The explanation tied the rule to the DTO layer that already exists here, which made the change feel justified.
}

%\textit{``The explanation tied the rule to the DTO layer that already exists here, which made the change feel justified.''}

P2 noted: 
\pquote{P2}{with DesignFix}{The explanation matched the factory pattern in this codebase, not just generic advice.}

%\textit{``The explanation matched the factory pattern in this codebase, not just generic advice.''} 

%P9, who struggled with the command hierarchy in the without DesignFix condition, found the explanation clarifying:
%\pquote{P9}{with DesignFix}{The explanation clarified the factory-to-subclass mapping, which I was shaky on.}

%\textit{``The explanation clarified the factory-to-subclass mapping, which I was shaky on.''}

Even when the suggested code was incorrect, the explanation's correct identification of key design concepts (such as 'command indirection' and 'caller/callee') enabled participants to find the right fix pattern independently. 
%- a phenomenon explored further in Theme 3.


\textbf{Theme 3: Verification Affordances Enable Appropriate Trust Calibration.} We observed that successful tool use depends on verification support, not just suggestion accuracy. For example, P14 received an \textit{incorrect} fix 
%(direct artifact call violating the rule) 
but a \textit{correct} explanation. Through the review step, P14 detected the inconsistency:

\pquote{P14}{with DesignFix}{The fix didn't match the explanation, but the keywords were right.}

%\textit{``The fix didn't match the explanation, but the keywords were right.''} 

\noindent Using explanation keywords, P14 found correct patterns and succeeded. P8 received a \textit{correct} fix but skipped verification, pasted incorrectly, and failed.
This suggests that the tool provided \textit{trust calibration mechanisms}: features enabling appropriate trust or distrust of AI suggestions. Less experienced participants valued verification features as safety nets:
\pquote{P13}{with DesignFix}{Having that intermediate step where I could see the code before it became part of the project gave me the confidence to proceed.}

%\textit{``Having that intermediate step where I could see the code before it became part of the project gave me the confidence to proceed.''}

\textbf{Theme 4: Cognitive Load Reduction.} 12 of 14 participants explicitly described reduced working memory and attention demands. The without DesignFix condition required simultaneously tracking file locations, understanding patterns, remembering previous exploration paths, and reasoning about the fix. P5 noted:
\pquote{P5}{with DesignFix}{Without the tool, I had to hold so much in my head: where files are, what the pattern looks like, what I've already tried. With the tool, I could just focus on understanding the fix rather than discovering it.}

%\textit{``Without the tool, I had to hold so much in my head: where files are, what the pattern looks like, what I've already tried. With the tool, I could just focus on understanding the fix rather than discovering it.''}

P9, who reached the time limit in the without DesignFix condition after ``getting lost jumping between the Command files,'' found the split view particularly helpful: 
\pquote{P9}{with DesignFix}{
It automatically opened the Factory class on the left and the fix on the right. I didn't have to search for the file or worry about losing my place, I could verify the AddCallee signature side-by-side.
}


%\textit{``It automatically opened the Factory class on the left and the fix on the right. I didn't have to search for the file or worry about losing my place - I could verify the AddCallee signature side-by-side.''}

\textbf{Theme 5: Workflow Integration and Friction Points.} Participants had varied preferences for fix application, revealing a tension between speed and safety (mentioned by 10 of 14 participants). Some appreciated the review step: P1 noted, 
\pquote{P1}{with DesignFix}{I actually liked that it opened in a new untitled tab first. Having the explanation text right above the code block helped me verify the fix before committing.}

%\textit{``I actually liked that it opened in a new untitled tab first. Having the explanation text right above the code block helped me verify the fix before committing.''} 

P5 used \textit{``the untitled file as a sandbox''} to make minor adjustments before accepting. Additionally, some participants found friction in the multi-step workflow, noting that clicking 'Accept' in the browser and then switching context to the IDE required too many steps.
%\textit{``It felt like too many clicks. I clicked `Accept' in the browser, then I had to switch context to the IDE.''} 

P10 experienced state management issues: 
\pquote{P10}{with DesignFix}{I clicked `Suggest Fix' twice and got two different suggestions... suddenly I had multiple `Untitled' tabs open. It became messy very quickly.}

%\textit{``I clicked `Suggest Fix' twice and got two different suggestions... suddenly I had multiple `Untitled' tabs open. It became messy very quickly.''} 

Nine participants requested multi-violation scanning to identify all related violations at once. Seven participants suggested surfacing related patterns from the codebase to enable verification against existing correct examples.

%\noindent Other participants suggested workflow improvements such as keyboard shortcuts to accept changes to reduce context switching.

\subsection{Failure Case Analysis}
\label{sec:failure-analysis}
% Add subsection text here.
Two cases merit detailed analysis for the insights they provide about system limitations and design implications.

\textbf{Case 1: Incorrect Fix Placement.} P8 represents the only failure in the DesignFix condition. The participant received a correct suggested fix 
%for the command-only artifact creation rule 
but pasted it in an incorrect location without careful verification. 
%P8 gave up the unassisted task at 12 minutes rather than continue searching. With DesignFix, P8 completed in 8 minutes but produced an incorrect result.
P8's workflow feedback reveals the issue: 
\pquote{P8}{with DesignFix}{This was much faster than my usual workflow. I didn't have to read any extra fluff... I saw the one-line constructor change, compared it in one second, and hit accept.}

%\textit{``This was much faster than my usual workflow. I didn't have to read any extra fluff... I saw the one-line constructor change, compared it in one second, and hit accept.''} 

\noindent The speed-oriented approach caused P8 to skip verification. This illustrates that while DesignFix dramatically reduces time and effort, developers must still engage in a review process.

\textbf{Case 2: Incorrect LLM Suggestion Recovered via Explanation.} P14 encountered a case where DesignFix's suggested code was incorrect --- it used a direct call that violated the command indirection rule ---  but the explanation correctly described the rule's requirements.

P14 detected the inconsistency: 
\pquote{P14}{with DesignFix}{
The tool actually gave me a bad patch - it suggested a direct call which was a violation. Because the workflow opened an untitled file first, I was able to catch this.
}

%\textit{``The tool actually gave me a bad patch - it suggested a direct call which was a violation. Because the workflow opened an untitled file first, I was able to catch this.''}

\noindent P14 used explanation keywords (``command indirection,'' ``caller/callee'') to find the correct pattern in FunctionCommand.java, then manually corrected the suggestion before accepting. This demonstrates that explanations serve as an independent verification mechanism: even when code is wrong, correct explanations enable skilled developers to recover.

\section{Threats to Validity}
\label{sec:threats}
% Add section text here.
\textit{Study design, task selection, and generalizability.}
The within-subjects design reduces between-participant variation but
introduces possible order, learning, and carryover effects. Condition
order was counterbalanced, but rule assignment was not fully
counterbalanced: some rules occurred only in one condition and others
were unevenly distributed. Consequently, condition effects may be
partly confounded with rule difficulty.

The tasks were purposively selected by both authors using
Evaluation~I outcomes and author-assessed difficulty criteria rather
than randomly sampled from the 15 CrowdCode rules. Conditioning task
selection on observed repair outcomes can influence the estimated
effect, while subjective task pairing may leave residual difficulty
differences between conditions. The reported results should therefore
be interpreted as evidence for these eight selected tasks rather than
as an unbiased estimate across all CrowdCode rules.

Three prompt-development violations were also retained in the
45-violation Evaluation~I dataset, which may introduce optimistic bias.
We report held-out results separately to make this effect visible. 

\textit{Tool-access and condition-design bias.}
We chose the checker-centered comparator because it represents the
existing workflow that DesignFix augments. Participants received the
rule, violating snippet, file path, click-to-open violation
localization, source repository, tutorial, and optional browser-based
tools. DesignFix added an integrated bundle of same-project examples,
focused context, explanation, repair-site guidance, proposed code, and
review support.

Accordingly, the study estimates the difference between these two
workflows. It neither isolates the contribution of individual
DesignFix features nor compares LLM access with no LLM access.
Browser-LLM use in the checker-centered condition was self-selected,
creating additional heterogeneity that we report descriptively rather
than interpreting causally.

\textit{Tutorial equivalence and carryover.}
All participants were taught both workflows before either task. This
equalized initial instruction but may have introduced cross-condition
contamination. Completing the first task could also provide repository
knowledge that carried into the second task; counterbalancing condition
order mitigates but does not eliminate this possibility.

\textit{Rule-language expressiveness.} DesignFix inherits the
expressiveness limits of its rule checker. srcML encodes only
syntactic information obtained from parsing: it operates on local
document information without a symbol table and performs no name
resolution, type binding, or control- or data-flow
analysis~\cite{collard2013srcml}. XPath queries over this
representation match syntactic patterns; they cannot express
constraints over statically computed semantic properties such as a
resolved call graph, dataflow facts, or protocol and ordering
properties~\cite{mehrpour2020rulepad}. Design rules requiring such
analyses are outside the rule language and were not part of our
evaluation.

\textit{Single-rater assessment.}
Explanation Correctness and Surrounding Code Pass Rate in Evaluation~I,
and repair correctness in the user study, were assessed by the first
author using predefined rubrics. We did not measure inter-rater
reliability, which is a threat to construct validity
~\cite{mcdonald2019reliability}.

Checker compliance and user-study compilation were tool-checked, but
determining explanation correctness and whether a modification was
unrelated or unnecessary required human judgment. Explicit condition
labels were withheld while study repairs were assessed; however,
masking may have been incomplete because task identity and patch
content could provide condition cues. A second rater might classify
borderline cases differently. We release the rubrics, outputs, diffs,
and labels so that these judgments can be independently audited.


\section{Discussion}
\label{sec:discussion}
% Add section text here.
On the selected tasks, participants had shorter recorded task durations with the integrated DesignFix workflow. 
%\subsection{The Navigation Burden in Unfamiliar Codebases}
This difference reflects the combined support provided by project-specific context, repair-site identification, explanation, proposed code, and review; not merely access to an LLM. Participants in the without DesignFix condition spent more than twice as long before making their first code modification (Theme~1, Section~\ref{sec:user-study-qualitative}), and several edited entirely wrong files without realizing their error. This suggests that design rule violation repair is primarily a \textit{navigation and comprehension} problem rather than a coding problem. DesignFix's primary value may therefore lie in codebase navigation and pattern identification rather than code generation.

%\subsection{Experience Stratification in Tool Adoption}

Developer experience shaped not just performance but workflow adaptation. Less experienced participants used the review workflow as intended, while experts like P11 bypassed the acceptance buttons yet still benefited substantially (Theme~1). This suggests LLM-assisted tools should support multiple interaction modes: guided workflows for novices and minimal-friction access to underlying reasoning for experts.

%\subsection{Explanation-Code Alignment as a Safety Mechanism}

The P14 case, where DesignFix provided incorrect code but a correct explanation, reveals an important safety property. Because explanations were prominently displayed alongside suggestions, P14 detected the inconsistency and recovered using explanation keywords. This suggests that tools should not sacrifice explanation quality for code accuracy, as explanations serve as an independent verification mechanism.

\section{Conclusions}
\label{sec:conclusion}
% Add section text here.
We presented DesignFix, an IDE-integrated workflow that transforms
checker-reported design-rule violations into project-grounded
explanations and reviewable repairs. Across 9,450 trials, the strongest
tested workflow achieved up to 90\% repair success, 89\% explanation
correctness, and 90\% preservation of unrelated code. In a
14-participant within-subjects study, DesignFix reduced mean task
duration by 36\%, increased correct repairs from 43\% to 93\%, and
significantly reduced five of six NASA-TLX dimensions. Together, these
results show that effective LLM support for design-rule repair extends
beyond generating code: connecting a violation to the relevant project
context, repair site, explanation, and review workflow can make
unfamiliar design constraints substantially easier to understand and
act upon. Future work should investigate multi-file repairs, longer-term
use, and generalization across additional codebases and LLM providers.

\section{Data Availability Statement}
\label{sec:data-availability}
% Add section text here.
The data that supports the findings of this study are available in the ``Supplementary material'' section of this submission. 
\bibliographystyle{ACM-Reference-Format}
\bibliography{references}

\end{document}